\documentclass{article}
\usepackage{pseudocode}

\begin{document}


\begin{algorithme}
\type{TableDeCodage}{\typePointeur{Noeud}}
\begin{enregistrement}{Noeud}
    \champEnregistrement{code}{CodeBinaire}
    \champEnregistrement{car}{\caractere}
    \champEnregistrement{listeSuivante}{TableDeCodage}

\end{enregistrement}
\end{algorithme}


\section{TableDeCodage}

\begin{algorithme}
    \fonction{obtenirTableDeCodage}{a : Arbre Huffman}{TableDeCodage}
    {non estVide(a)}
    {res : TableDeCodage}
    
    début
        res.car < obtenirCaractere(a.statistiques)
        res.code < obtenirCodeBinaire(res.car)
        listeSuivante < NIL
        retourner res
    fin

    {\affecter{file}{creerFileDePriorite(fic)}
    \affecter{arbreHuffman}{creerArbreHuffman(file)}
    \affecter{table}{creerTableDeCodage(arbreHuffman)}}
\end{algorithme}



\begin{algorithme}
    \fonction{obtenirCode}{tab: TableDeCodage, car : \caractere}{CodeBinaire}
    {estPresent(tab,car)}
    {}
    début
        si non(estVide(tab)) alors 
            si tab.car=car alors
                retourner tab.code
            sinon
                retourner obtenirCode(tab.listeSuivante)
            finsi
        finsi
    fin
    {}
\end{algorithme}

\begin{algorithme}
    \fonction{obtenirCarac}{tab: TableDeCodage, code : CodeBinaire}{\caractere}
    {}
    {}
    début
        si non(estVide(tab)) alors 
            si tab.code=code alors
                retourner tab.car
            sinon
                retourner obtenirCode(tab.listeSuivante)
            finsi
        finsi
    fin
    {}
\end{algorithme}

\begin{algorithme}
    \procedure{ajouter}{\paramEntreeSortie tab : TableDeCodage, \paramEntree car : \caractere, code : CodeBinaire}
    {}
    {temp : TableDeCodage}
    début
        temp < tab
        allouer(tab)
        tab.car < car
        tab.code < code
        tab.listeSuivante < temp
    fin
    {}
\end{algorithme}

\begin{algorithme}
    \procedure{supprimer}{\paramEntreeSortie tab : TableDeCodage, \paramEntree car : \caractere, code : CodeBinaire}
    {}
    {}

    {}
\end{algorithme}

\begin{algorithme}
    \procedure{estPresent}{\paramEntreeSortie tab : TableDeCodage, \paramEntree car : \caractere, code : CodeBinaire}
    {non(estVide(tab))}
    {}

    {}
\end{algorithme}

\end{document}
